% Header: General study / work / project template
%%%%%%%%%%%%%%%%%%%%%%%%%%%%%%%%%%%%%%%%%%%%%%%%%%%%%%%%%%%%%%%%%%%

\documentclass[11pt,a4paper]{scrartcl}

% Layout
\usepackage[margin=2.5cm]{geometry}
\usepackage[onehalfspacing]{setspace}

% General packages
\usepackage{graphicx}
\usepackage{enumitem}
\usepackage{tcolorbox}
\tcbuselibrary{breakable}
\usepackage[breaklinks=true,colorlinks=true,linkcolor=blue,urlcolor=blue,citecolor=blue]{hyperref}
\usepackage{amsmath,amsthm,amssymb,mathtools}
\usepackage{icomma}
\usepackage[T1]{fontenc}
\usepackage[utf8]{inputenc}

% Language: choose one
% \usepackage[ngerman]{babel}
\usepackage[english]{babel}

% Units / tables / floats / references
\usepackage[locale = UK,space-before-unit=true,per-mode = symbol]{siunitx}
\usepackage{booktabs,multirow}
\usepackage{placeins}
\usepackage[natbib,abbreviate=true,doi=false,style=numeric-comp,giveninits=true,sorting=none]{biblatex}
\usepackage{csquotes}

% Optional bibliography
% \addbibresource{references.bib}

% Graphics paths
\graphicspath{{Bilder/} {images/} {figures/} {chapters/}}

% Optional math shortcuts
\newcommand{\R}{\mathbb{R}}
\newcommand{\N}{\mathbb{N}}
\newcommand{\Q}{\mathbb{Q}}
\newcommand{\set}[1]{\left\{ #1 \right\}}
\newcommand{\abs}[1]{\left| #1 \right|}
\newcommand{\norm}[1]{\left\lVert #1 \right\rVert}
\newcommand{\ol}[1]{\overline{#1}}

% Optional units
\DeclareSIUnit{\dBm}{dBm}

% Box styles
\newtcolorbox{calculationbox}[1][]{
    title=Calculation,
    breakable,
    #1
}

\newtcolorbox{solutionbox}[1][]{
    title=Solution,
    breakable,
    #1
}

\newtcolorbox{answerbox}[1][]{
    title=Answer,
    breakable,
    #1
}

\newtcolorbox{notebox}[1][]{
    title=Notes,
    breakable,
    #1
}

\newtcolorbox{sidecalcbox}[1][]{
    title=Side Calculation,
    breakable,
    #1
}

%%%%%%%%%%%%%%%%%%%%%%%%%%%%%%%%%%%%%%%%%%%%%%%%%%%%%%%%%%%%%%%%%%%
\begin{document}

\title{Physik 4 - Blatt 04}
\author{Tobias Laser}
\date{\today}
\maketitle
\vfill
\thispagestyle{empty}

\pagenumbering{arabic}

\paragraph{Aufgabe 1: Bethe-Weizsäcker: der Coulombterm}

Für schwere Atomkerne können Sie diese als homogen geladene Kugel betrachten.
Dieses Modell kann man zur Herleitung des Coulombterms aus der Bethe-Weizsäcker-Formel benutzen.

\begin{enumerate}[label=(\alph*)]

    \item Wie groß ist die elektrostatische Abstoßungsenergie $E_{\mathrm{ab}}$ einer homogen geladenen Kugel mit Gesamtladung $Q$ und Radius $R$? Nehmen Sie dazu Ihre Aufzeichnungen aus der ersten Anwesenheitsübung her.

    \begin{solutionbox}
        Gesucht ist die elektrostatische Selbstenergie einer homogen geladenen Kugel.
        Diese erhält man, indem man sich vorstellt, dass die Kugel schrittweise aus unendlich dünnen Kugelschalen aufgebaut wird.

        Wenn die Kugel gerade den Radius $r$ besitzt, dann ist die bis dahin eingebrachte Ladung wegen der homogenen Ladungsverteilung proportional zum Volumen:
        \[
            q(r)=Q\frac{r^3}{R^3}.
        \]

        Das Potential an der Oberfläche einer Kugel mit Radius $r$ und Ladung $q(r)$ ist
        \[
            \phi(r)=\frac{1}{4\pi\varepsilon_0}\frac{q(r)}{r}.
        \]

        Fügt man nun eine infinitesimale zusätzliche Ladung $dq$ auf dieser Oberfläche hinzu, so muss dafür die Arbeit
        \[
            dE=\phi(r)\,dq 
        \]
        aufgebracht werden.n
        Wegen
        \[
            q(r)=Q\frac{r^3}{R^3}
        \]
        folgt durch Differentiation nach $r$:
        \[
            dq=\frac{3Q}{R^3}r^2\,dr.
        \]

        Damit wird
        \[
            dE
            =\frac{1}{4\pi\varepsilon_0}\frac{q(r)}{r}\,dq
            =\frac{1}{4\pi\varepsilon_0}\frac{Q r^3/R^3}{r}\cdot \frac{3Q}{R^3}r^2\,dr.
        \]

        Dies vereinfacht sich zu
        \[
            dE
            =\frac{1}{4\pi\varepsilon_0}\frac{3Q^2}{R^6}r^4\,dr.
        \]

        Nun integrieren wir von $r=0$ bis $r=R$:
        \[
            E_{\mathrm{ab}}
            =\int_0^R dE
            =\frac{3Q^2}{4\pi\varepsilon_0 R^6}\int_0^R r^4\,dr.
        \]

        Mit
        \[
            \int_0^R r^4\,dr=\frac{R^5}{5}
        \]
        erhält man
        \[
            E_{\mathrm{ab}}
            =\frac{3Q^2}{4\pi\varepsilon_0 R^6}\cdot \frac{R^5}{5}
            =\frac{3}{5}\frac{1}{4\pi\varepsilon_0}\frac{Q^2}{R}.
        \]
    \end{solutionbox}

    \item Setzen Sie nun diese Energie $E_{\mathrm{ab}}$ gleich dem Coulombterm und verwenden Sie die Formel für den Rutherford-Radius
    \[
        R=R_0A^{1/3},
    \]
    um daraus eine Formel für $a_C$ herzuleiten.

    \begin{solutionbox}
        In einem Atomkern mit Protonenzahl $Z$ ist die Gesamtladung
        \[
            Q=Ze.
        \]

        Setzt man dies in die in Teil (a) gefundene Formel ein, so erhält man
        \[
            E_{\mathrm{ab}}
            =\frac{3}{5}\frac{1}{4\pi\varepsilon_0}\frac{(Ze)^2}{R}
            =\frac{3}{5}\frac{e^2}{4\pi\varepsilon_0}\frac{Z^2}{R}.
        \]

        Nun verwenden wir den Rutherford-Radius
        \[
            R=R_0A^{1/3}.
        \]
        Daraus folgt
        \[
            E_{\mathrm{ab}}
            =\frac{3}{5}\frac{e^2}{4\pi\varepsilon_0}\frac{Z^2}{R_0A^{1/3}}
            =\left(\frac{3}{5}\frac{e^2}{4\pi\varepsilon_0R_0}\right)\frac{Z^2}{A^{1/3}}.
        \]

        Der Coulombterm der Bethe-Weizsäcker-Formel hat genau die Struktur
        \[
            E_C=a_C\frac{Z^2}{A^{1/3}}.
        \]

        Durch Koeffizientenvergleich erhält man direkt
        \[
            \boxed{
            a_C=\frac{3}{5}\frac{e^2}{4\pi\varepsilon_0R_0}
            }.
        \]
    \end{solutionbox}

    \item Welcher Wert ergibt sich für $R_0=1.2\,\mathrm{fm}$ und welcher für $R_0=1.3\,\mathrm{fm}$? Für wie gut halten Sie also diese Modellvorstellung?

    \begin{solutionbox}
        Wir benutzen den bekannten numerischen Wert
        \[
            \frac{e^2}{4\pi\varepsilon_0}\approx 1.44\,\mathrm{MeV\,fm}.
        \]

        Damit wird
        \[
            a_C=\frac{3}{5}\frac{1.44\,\mathrm{MeV\,fm}}{R_0}.
        \]

        \textbf{Fall 1:} $R_0=1.2\,\mathrm{fm}$

        Dann folgt
        \[
            a_C=\frac{3}{5}\frac{1.44}{1.2}\,\mathrm{MeV}
            =\frac{3}{5}\cdot 1.2\,\mathrm{MeV}
            =0.72\,\mathrm{MeV}.
        \]

        \textbf{Fall 2:} $R_0=1.3\,\mathrm{fm}$

        Dann folgt
        \[
            a_C=\frac{3}{5}\frac{1.44}{1.3}\,\mathrm{MeV}
            \approx \frac{3}{5}\cdot 1.1077\,\mathrm{MeV}
            \approx 0.6646\,\mathrm{MeV}.
        \]

        Also erhält man Werte im Bereich von ungefähr
        \[
            a_C\approx 0.66\text{ bis }0.72\,\mathrm{MeV}.
        \]

        Der in der Bethe-Weizsäcker-Formel üblicherweise verwendete Wert liegt etwa bei
        \[
            a_C\approx 0.71\,\mathrm{MeV}.
        \]

        Das einfache Modell der homogen geladenen Kugel liefert also bereits eine sehr gute Größenordnung und sogar numerisch recht passende Werte.
    \end{solutionbox}

\end{enumerate}

\paragraph{Aufgabe 2: Bethe-Weizsäcker: der Volumenterm}

Der Volumenterm der Bethe-Weizsäcker-Formel lautet $a_VA$.
Das soll nun an einem einfachen Modell nachvollzogen werden.
Gehen Sie dazu davon aus, dass aufgrund der geringen Reichweite der Kernkraft nur die nächsten Nachbarn miteinander wechselwirken.
Dabei soll bei der Bindung zweier Nukleonen eine Wechselwirkungsenergie $\varepsilon$ frei werden.

\begin{enumerate}[label=(\alph*)]

    \item Nehmen Sie nun an, die Verteilung der Nukleonen im Kern entspricht der dichtesten Kugelpackung. Wieviele nächste Nachbarn $N_{\mathrm{NN}}$ besitzt dann ein Nukleon?

    \begin{solutionbox}
        In der dichtesten Kugelpackung ist jede Kugel von möglichst vielen anderen Kugeln direkt umgeben.
        Ein bekanntes Resultat aus der Geometrie der dichtesten Kugelpackung ist, dass jede Kugel genau 12 nächste Nachbarn besitzt.

        Übertragen auf den Kern bedeutet dies:
        \[
            N_{\mathrm{NN}}=12.
        \]

        Da laut Aufgabenstellung Oberflächeneffekte ignoriert werden dürfen, gilt diese Zahl für alle Nukleonen im Modell.
    \end{solutionbox}

    \item Da ein Kern aus $A$ Nukleonen besteht, würde man erwarten, dass sich die gesamte hier diskutierte Wechselwirkungsenergie über
    \[
        E_{\mathrm{WW}}=k\varepsilon A
    \]
    berechnet. Wie groß sollte also im Sinne von Teil (a) die Konstante $k$ sein?

    \begin{solutionbox}
        Nach Teil (a) hat jedes Nukleon 12 nächste Nachbarn.
        Würde man die Zahl der Bindungen einfach als
        \[
            12A
        \]
        ansetzen, dann würde jede Bindung doppelt gezählt:
        einmal vom ersten Nukleon aus und einmal vom zweiten Nukleon aus.

        Deshalb muss man durch 2 teilen. Die tatsächliche Zahl der Nachbarbindungen ist also
        \[
            \frac{12A}{2}=6A.
        \]

        Jede dieser Bindungen liefert die Energie $\varepsilon$.
        Also ist die gesamte Wechselwirkungsenergie
        \[
            E_{\mathrm{WW}}=6A\,\varepsilon.
        \]

        Vergleicht man dies mit
        \[
            E_{\mathrm{WW}}=k\varepsilon A,
        \]
        dann folgt unmittelbar
        \[
            \boxed{k=6}.
        \]
    \end{solutionbox}

    \item Nun soll $\varepsilon$ anhand von Daten bestimmt werden. Um einen möglichst einfachen Kern zu betrachten, nutzen wir hier die Bindungsenergie von ${}^2\mathrm{H}$ (Deuterium), die $2.23\,\mathrm{MeV}$ beträgt. Was ergibt sich in diesem Fall für $\varepsilon$? Welcher Wert folgt somit für $a_V$ und wie vergleicht sich das mit dem bisher benutzten Wert von $a_V=15.85\,\mathrm{MeV}$?

    \begin{solutionbox}
        Das Deuterium ${}^2\mathrm{H}$ besteht aus genau zwei Nukleonen:
        einem Proton und einem Neutron.
        Zwischen diesen beiden Nukleonen gibt es genau eine Bindung.

        In diesem einfachen Modell ist die gesamte Bindungsenergie des Deuteriums daher direkt gleich der Bindungsenergie einer einzelnen Nachbarbindung:
        \[
            \varepsilon = 2.23\,\mathrm{MeV}.
        \]

        Aus Teil (b) wissen wir, dass im Modell für große Kerne
        \[
            E_{\mathrm{WW}}=6\varepsilon A
        \]
        gilt.
        Der Volumenterm der Bethe-Weizsäcker-Formel hat die Form
        \[
            a_VA.
        \]

        Daher identifiziert man
        \[
            a_V=6\varepsilon.
        \]

        Mit $\varepsilon=2.23\,\mathrm{MeV}$ folgt
        \[
            a_V=6\cdot 2.23\,\mathrm{MeV}=13.38\,\mathrm{MeV}.
        \]

        Verglichen mit dem üblichen Wert
        \[
            a_V=15.85\,\mathrm{MeV}
        \]
        ist das Modellresultat kleiner.
        Die Differenz beträgt
        \[
            15.85-13.38=2.47\,\mathrm{MeV}.
        \]

        Relativ gesehen ergibt sich
        \[
            \frac{2.47}{15.85}\approx 0.156 \approx 15.6\%.
        \]

        Das Modell erklärt also die Größenordnung des Volumenterms recht gut, ist aber nicht exakt.
        Das ist plausibel, weil das Modell sehr stark vereinfacht:
        Es berücksichtigt z.\,B. nicht die komplizierte Struktur der Kernkraft und keine feineren quantenmechanischen Effekte.
    \end{solutionbox}

\end{enumerate}

\paragraph{Aufgabe 3: Stabilstes Nuklid einer Isobare}

Verwenden Sie im Folgenden die Bethe-Weizsäcker-Formel:
\[
M(A, Z) = (A - Z)m_N + Zm_P + Zm_e
-
\left(
a_VA - a_SA^{2/3} - a_CZ^2A^{-1/3} - a_A\frac{(A - 2Z)^2}{A} + \delta a_P A^{-1/2}
\right)c^{-2},
\]
mit den Koeffizienten
\[
a_V = 15.58\,\mathrm{MeV}, \qquad
a_S = 16.91\,\mathrm{MeV}, \qquad
a_C = 0.71\,\mathrm{MeV}, \qquad
a_A = 23.21\,\mathrm{MeV}, \qquad
a_P = 11.46\,\mathrm{MeV}.
\]

Im Folgenden wird die Massenabhängigkeit $M(A,Z)$ einer Isobarenreihe ($A=\text{const.}$) betrachtet.

\begin{enumerate}[label=(\alph*)]

    \item Berechnen Sie mit Hilfe der Bethe-Weizsäcker-Formel für eine konstante Nukleonenzahl die Kernladungszahl $Z_0$, bei der die Kernmasse minimal wird. Nehmen Sie hierfür $\delta=0$ an.

    \begin{solutionbox}
        Gesucht ist das Minimum der Masse $M(A,Z)$ für festes $A$.
        Da $A$ konstant ist, sind alle Terme, die nur von $A$ abhängen, für die Minimierung irrelevant.
        Außerdem setzen wir gemäß Aufgabenstellung
        \[
            \delta=0.
        \]

        Wir betrachten also nur die $Z$-abhängigen Terme.
        Zweckmäßig multiplizieren wir die Massenformel zunächst mit $c^2$, damit wir direkt mit Energien arbeiten können:
        \[
            M(A,Z)c^2
            =(A-Z)m_Nc^2+Z(m_P+m_e)c^2
            -a_VA+a_SA^{2/3}+a_CZ^2A^{-1/3}+a_A\frac{(A-2Z)^2}{A}.
        \]

        Für die Ableitung nach $Z$ können wir die konstanten Terme weglassen.
        Es bleibt
        \[
            f(Z)= -Zm_Nc^2 + Z(m_P+m_e)c^2 + a_CZ^2A^{-1/3} + a_A\frac{(A-2Z)^2}{A}.
        \]

        Zusammengefasst:
        \[
            f(Z)= Z(m_P+m_e-m_N)c^2 + a_CZ^2A^{-1/3} + a_A\frac{(A-2Z)^2}{A}.
        \]

        Nun leiten wir nach $Z$ ab.

        Für den linearen Term gilt:
        \[
            \frac{d}{dZ}\left[Z(m_P+m_e-m_N)c^2\right]
            =(m_P+m_e-m_N)c^2.
        \]

        Für den Coulombterm:
        \[
            \frac{d}{dZ}\left[a_CZ^2A^{-1/3}\right]
            =2a_CZA^{-1/3}.
        \]

        Für den Asymmetrieterm:
        \[
            \frac{d}{dZ}\left[a_A\frac{(A-2Z)^2}{A}\right]
            =a_A\frac{2(A-2Z)(-2)}{A}
            =-4a_A\frac{A-2Z}{A}.
        \]

        Damit ist insgesamt
        \[
            \frac{df}{dZ}
            =(m_P+m_e-m_N)c^2 + 2a_CZA^{-1/3} -4a_A\frac{A-2Z}{A}.
        \]

        Für ein Extremum setzen wir
        \[
            \frac{df}{dZ}=0.
        \]
        Also
        \[
            (m_P+m_e-m_N)c^2 + 2a_CZA^{-1/3} -4a_A\frac{A-2Z}{A}=0.
        \]

        Den letzten Term schreiben wir aus:
        \[
            -4a_A\frac{A-2Z}{A}
            =-4a_A+\frac{8a_A}{A}Z.
        \]

        Somit wird die Gleichung
        \[
            (m_P+m_e-m_N)c^2 + 2a_CZA^{-1/3} -4a_A + \frac{8a_A}{A}Z = 0.
        \]

        Wir sammeln alle $Z$-Terme:
        \[
            Z\left(2a_CA^{-1/3}+\frac{8a_A}{A}\right)
            =4a_A-(m_P+m_e-m_N)c^2.
        \]

        Damit folgt
        \[
            \boxed{
            Z_0=
            \frac{4a_A-(m_P+m_e-m_N)c^2}
            {2a_CA^{-1/3}+\dfrac{8a_A}{A}}
            }.
        \]

        Multipliziert man Zähler und Nenner noch mit $A$, so erhält man die äquivalente Form
        \[
            \boxed{
            Z_0=
            \frac{A\left[4a_A-(m_P+m_e-m_N)c^2\right]}
            {2a_CA^{2/3}+8a_A}
            }.
        \]
    \end{solutionbox}

    \item Berechnen Sie $Z_0$ für $A=25$ und für $A=67$ und vergleichen Sie Ihr Ergebnis mit der Nuklidkarte. Diskutieren Sie das Ergebnis. Zeichnen Sie in beiden Fällen eine Grafik der Massenparabeln im Bereich $Z_0\pm 3$, wobei Sie sowohl die Masse laut Massenformel als auch die gemessenen Massen eintragen.

    \begin{solutionbox}
        Zunächst benötigen wir die Massendifferenz
        \[
            (m_P+m_e-m_N)c^2.
        \]
        Mit den bekannten Ruheenergien gilt näherungsweise
        \[
            m_Pc^2\approx 938.272\,\mathrm{MeV},\qquad
            m_ec^2\approx 0.511\,\mathrm{MeV},\qquad
            m_Nc^2\approx 939.565\,\mathrm{MeV}.
        \]
        Also
        \[
            (m_P+m_e-m_N)c^2
            \approx 938.272+0.511-939.565
            \approx -0.782\,\mathrm{MeV}.
        \]

        Damit wird die Formel aus Teil (a)
        \[
            Z_0=
            \frac{4a_A-(m_P+m_e-m_N)c^2}
            {2a_CA^{-1/3}+\dfrac{8a_A}{A}}
        \]
        zu
        \[
            Z_0=
            \frac{4a_A+0.782}{2a_CA^{-1/3}+\dfrac{8a_A}{A}}.
        \]

        Da
        \[
            a_A=23.21\,\mathrm{MeV},
        \]
        ist
        \[
            4a_A+0.782=4\cdot 23.21+0.782=93.622\,\mathrm{MeV}.
        \]

        \textbf{Fall 1: $A=25$}

        Zuerst berechnen wir
        \[
            A^{-1/3}=25^{-1/3}\approx 0.342.
        \]
        Dann
        \[
            2a_CA^{-1/3}=2\cdot 0.71\cdot 0.342\approx 0.486.
        \]
        Weiter
        \[
            \frac{8a_A}{A}=\frac{8\cdot 23.21}{25}=\frac{185.68}{25}\approx 7.427.
        \]
        Also ist der Nenner
        \[
            0.486+7.427=7.913.
        \]
        Damit folgt
        \[
            Z_0(25)=\frac{93.622}{7.913}\approx 11.83.
        \]

        \textbf{Fall 2: $A=67$}

        Nun
        \[
            A^{-1/3}=67^{-1/3}\approx 0.2465.
        \]
        Daher
        \[
            2a_CA^{-1/3}=2\cdot 0.71\cdot 0.2465\approx 0.350.
        \]
        Und
        \[
            \frac{8a_A}{A}=\frac{8\cdot 23.21}{67}=\frac{185.68}{67}\approx 2.772.
        \]
        Also ist der Nenner
        \[
            0.350+2.772=3.122.
        \]
        Damit ergibt sich
        \[
            Z_0(67)=\frac{93.622}{3.122}\approx 29.998 \approx 30.00.
        \]

        Da die Kernladungszahl ganzzahlig sein muss, erwartet man die stabilsten Isobare also bei
        \[
            Z=12 \quad \text{für } A=25,
        \]
        und
        \[
            Z=30 \quad \text{für } A=67.
        \]

        Diese Werte entsprechen
        \[
            {}^{25}\mathrm{Mg}\quad (Z=12)
        \]
        bzw.
        \[
            {}^{67}\mathrm{Zn}\quad (Z=30).
        \]

        Das stimmt gut mit der Nuklidkarte überein:
        Die Bethe-Weizsäcker-Formel sagt die Lage des Minimums der Isobarenparabel also gut voraus.

        Für die geforderte Zeichnung trägt man für die theoretische Parabel die Werte im Bereich
        \[
            Z_0\pm 3
        \]
        auf, also ungefähr

        für $A=25$:
        \[
            Z=9,10,11,12,13,14,15,
        \]
        und für $A=67$:
        \[
            Z=27,28,29,30,31,32,33.
        \]

        Dabei berechnet man jeweils $M(A,Z)$ aus der Bethe-Weizsäcker-Formel und trägt zusätzlich die gemessenen Massen aus der Nuklidkarte als einzelne Punkte ein.

        Qualitativ erwartet man:
        \begin{itemize}
            \item eine nach oben geöffnete Parabel als theoretische Kurve,
            \item ihr Minimum bei etwa $Z=12$ bzw. $Z=30$,
            \item kleine Abweichungen der realen Massen von der idealen Parabel, da Schaleneffekte und andere mikroskopische Strukturen in der einfachen Massenformel nicht vollständig beschrieben werden.
        \end{itemize}
    \end{solutionbox}

\end{enumerate}

\end{document}